\documentclass[a4paper,10pt]{article}\usepackage[]{graphicx}\usepackage[]{xcolor}
% maxwidth is the original width if it is less than linewidth
% otherwise use linewidth (to make sure the graphics do not exceed the margin)
\makeatletter
\def\maxwidth{ %
  \ifdim\Gin@nat@width>\linewidth
    \linewidth
  \else
    \Gin@nat@width
  \fi
}
\makeatother

\definecolor{fgcolor}{rgb}{0.345, 0.345, 0.345}
\newcommand{\hlnum}[1]{\textcolor[rgb]{0.686,0.059,0.569}{#1}}%
\newcommand{\hlsng}[1]{\textcolor[rgb]{0.192,0.494,0.8}{#1}}%
\newcommand{\hlcom}[1]{\textcolor[rgb]{0.678,0.584,0.686}{\textit{#1}}}%
\newcommand{\hlopt}[1]{\textcolor[rgb]{0,0,0}{#1}}%
\newcommand{\hldef}[1]{\textcolor[rgb]{0.345,0.345,0.345}{#1}}%
\newcommand{\hlkwa}[1]{\textcolor[rgb]{0.161,0.373,0.58}{\textbf{#1}}}%
\newcommand{\hlkwb}[1]{\textcolor[rgb]{0.69,0.353,0.396}{#1}}%
\newcommand{\hlkwc}[1]{\textcolor[rgb]{0.333,0.667,0.333}{#1}}%
\newcommand{\hlkwd}[1]{\textcolor[rgb]{0.737,0.353,0.396}{\textbf{#1}}}%
\let\hlipl\hlkwb

\usepackage{framed}
\makeatletter
\newenvironment{kframe}{%
 \def\at@end@of@kframe{}%
 \ifinner\ifhmode%
  \def\at@end@of@kframe{\end{minipage}}%
  \begin{minipage}{\columnwidth}%
 \fi\fi%
 \def\FrameCommand##1{\hskip\@totalleftmargin \hskip-\fboxsep
 \colorbox{shadecolor}{##1}\hskip-\fboxsep
     % There is no \\@totalrightmargin, so:
     \hskip-\linewidth \hskip-\@totalleftmargin \hskip\columnwidth}%
 \MakeFramed {\advance\hsize-\width
   \@totalleftmargin\z@ \linewidth\hsize
   \@setminipage}}%
 {\par\unskip\endMakeFramed%
 \at@end@of@kframe}
\makeatother

\definecolor{shadecolor}{rgb}{.97, .97, .97}
\definecolor{messagecolor}{rgb}{0, 0, 0}
\definecolor{warningcolor}{rgb}{1, 0, 1}
\definecolor{errorcolor}{rgb}{1, 0, 0}
\newenvironment{knitrout}{}{} % an empty environment to be redefined in TeX

\usepackage{alltt}
\usepackage[utf8]{inputenc}
\usepackage{amsmath}                       % for mathematical features 
\usepackage[left=2.5cm,top=3cm,bottom=3cm,right=2.5cm]{geometry}   % text margins
\usepackage{booktabs}                      % for booktabs in print(xtable)).
\usepackage{authblk}                       % for footnote style author/affiliation
\usepackage{amssymb}
\usepackage{amsthm}
% LaTeX package for captions customization:
\usepackage[margin=10pt,font=small,labelfont=bf,format=hang]{caption}
\usepackage{hyperref}                      % Hyperlinks and \url{}
\hypersetup{
    colorlinks=true, % Enable colored text for links instead of drawing boxes
    linkcolor=blue,  % Color for internal document links (e.g., sections, equations, figures, tables)
    citecolor=blue,  % Color for bibliographic citation links (e.g., \cite{})
    urlcolor=blue    % Color for external web URLs (e.g., \url{})
}
\IfFileExists{upquote.sty}{\usepackage{upquote}}{}
\begin{document}
\title{Statistics in Health Sciences - Assignment 1}
\author{G05: Judit Félelz Guerrero (17048330), Lídia Abio Rojo (1705251), Catalina Mascaró Català (1708159), Elia Campos Villaro (1703842)}
\date{28th October 2026}



\maketitle

\newpage

\section{Introduction}
The aim of this assignment is to apply the framework proposed by Basagaña et al. to the data set copd in the file copd redux.RData. Data copd is a subset of 45 variables of the original data analyzed in the cited pape.


\section{Analysis}

\subsection{Description of missingness}

Provide the same missingness description than in the original analysis. Specifically, reproduce the following text that is shown in the original paper:

\begin{quotation}
``The subset used here included \textcolor{red}{342} participants and \textcolor{red}{85} variables. The range of missing data in the different variables went from \textcolor{red}{0\%} to \textcolor{red}{47.7\%}, with \textcolor{red}{68.2\%} of the variables having less than 5\% of the data missing. Only \textcolor{red}{13.7\%} of the participants had complete data on all \textcolor{red}{85} variables. Overall, \textcolor{red}{5.9\%} of the information was missing; that is, \textcolor{red}{5.9\%} of the total of \textcolor{red}{29,070} cells (\textcolor{red}{$85 \times 342 = 29,070$}) had missing data.''
\end{quotation}

Instructions:

\begin{itemize}
    \item Note that, since you are working with a subset of the original data set, red numbers above are not expected to be the same than in your results.
    
    \item In your \textbf{\texttt{Rnw}} document, the text above must be able to automatically update all red numbers in the case of changing the values of \textbf{\texttt{copd}}. I.e., if data change, then all red numbers must also automatically update accordingly in your \textbf{\texttt{pdf}} when compiling your \textbf{\texttt{Rnw}}. As a hint, you can follow this procedure:
    
    \begin{enumerate}
        \item For each red number, create an \textsf{R} object, which will depend on data and gives the desired quantity (i.e., if data change, the value in the object changes accordingly),
        \item To force rounding a quantity to a given decimal precision, you can use \texttt{sprintf()} (as seen in Exercise 0). Do not use \texttt{round()},
        \item To get the thousands comma separator as in ``29,070'', you can use \texttt{prettyNum()},
        \item Once you have an \textsf{R} object for each red number formatted as desired, use \texttt{\textbackslash Sexpr\{\}} to insert each of them inline in \LaTeX{} (as seen in Exercise 0).
    \end{enumerate}
\end{itemize}









\begin{knitrout}\footnotesize
\definecolor{shadecolor}{rgb}{0.969, 0.969, 0.969}\color{fgcolor}\begin{kframe}
\begin{alltt}
\hldef{> }\hldef{n_obs} \hlkwb{<-} \hlkwd{nrow}\hldef{(copd)}
\hldef{> }\hldef{n_vars} \hlkwb{<-} \hlkwd{ncol}\hldef{(copd)}
\hldef{> }\hldef{nans} \hlkwb{<-} \hlnum{100}\hlopt{*}\hlkwd{colMeans}\hldef{(}\hlkwd{is.na}\hldef{(copd))}
\hldef{> }\hldef{min_nan} \hlkwb{<-} \hlkwd{sprintf}\hldef{(}\hlsng{"%.2g"}\hldef{,} \hlkwd{min}\hldef{(nans))} \hlcom{# MIRAR}
\hldef{> }\hldef{max_nan} \hlkwb{<-} \hlkwd{sprintf}\hldef{(}\hlsng{"%.2g"}\hldef{,} \hlkwd{max}\hldef{(nans))}
\hldef{> }\hldef{less_5} \hlkwb{<-} \hlkwd{sprintf}\hldef{(}\hlsng{"%.2g"}\hldef{,} \hlnum{100}\hlopt{*}\hlkwd{length}\hldef{(nans[nans}\hlopt{<}\hlnum{5}\hldef{])}\hlopt{/}\hldef{n_vars)}
\hldef{> }\hldef{comp_cases} \hlkwb{<-} \hlkwd{sprintf}\hldef{(}\hlsng{"%.2g"}\hldef{,} \hlnum{100}\hlopt{*}\hlkwd{mean}\hldef{(}\hlkwd{complete.cases}\hldef{(copd)))}
\hldef{> }\hldef{miss_total} \hlkwb{<-}  \hlkwd{sprintf}\hldef{(}\hlsng{"%.2g"}\hldef{,} \hlnum{100}\hlopt{*}\hlkwd{mean}\hldef{(}\hlkwd{is.na}\hldef{(copd)))}
\hldef{> }\hldef{total_cells} \hlkwb{=} \hlkwd{prettyNum}\hldef{(n_obs}\hlopt{*}\hldef{n_vars,} \hlkwc{big.mark} \hldef{=} \hlsng{","}\hldef{,} \hlkwc{decimal.mark} \hldef{=} \hlsng{"."}\hldef{)}
\end{alltt}
\end{kframe}
\end{knitrout}

“The subset used here included 342 participants and 45 variables. The range of missing data in the different variables went from 0\% to 32\%, with 69\% of the variables having less than 5\% of the data missing. Only 33\% of the participants had complete data on all 45 variables. Overall, 5.4\% of the information was missing; that is, 5.4\% of the total of 15,390 cells (45 × 342 = 15,390) had missing data.”

\end{document}
